\documentclass{article}

% NeurIPS 2026 style (mock include — sty file expected at compile time)
\usepackage[final]{neurips_2026}

\usepackage[utf8]{inputenc}
\usepackage[T1]{fontenc}
\usepackage{hyperref}
\usepackage{url}
\usepackage{booktabs}
\usepackage{amsfonts}
\usepackage{amsmath}
\usepackage{amssymb}
\usepackage{nicefrac}
\usepackage{microtype}
\usepackage{xcolor}
\usepackage{graphicx}

\title{Sparse Rationale-Constrained Attention for Hate Speech Detection}

\author{%
  Anonymous Authors\\
  Anonymous Institution\\
  \texttt{anonymous@example.org}
}

\begin{document}

\maketitle

\begin{abstract}
Hate speech classifiers deployed in moderation pipelines must justify their decisions, yet
attention weights are routinely shown to be unfaithful explanations of model behaviour.
Sparse Rationale-Constrained Attention (SRA), recently introduced as a remedy, replaces
softmax with sparsemax and supervises a single attention head with token-level rationales.
We revisit SRA on HateXplain and show that, while the structural insight is correct, the
\emph{instantiation} is severely under-constrained: a single supervised head leaks
explanatory mass into the remaining eleven through the residual stream. We propose an
all-layer, all-head sparsemax projection with mean-squared rationale supervision (C4) and
a top-six layer variant (C5). Across five seeds we observe (i) classification parity with
the softmax baseline (macro-F1 $0.6823 \pm 0.0063$ vs.\ $0.6814 \pm 0.0086$, TOST $p{=}0.047$
at $\pm 1.0$ pp), (ii) a $111\%$ relative gain in ERASER comprehensiveness AOPC ($0.338$
vs.\ $0.160$ for SRA, Wilcoxon $p{=}0.031$, $d{=}8.36$), and (iii) a $47\%$ reduction in
sufficiency AOPC, all at zero additional inference cost. We honestly disclose that the
pre-registered $2{\times}$ adversarial-swap robustness threshold (H4) is missed: the
observed C4/C2 KL ratio is $1.47{\times}$. A structural argument (Proposition~1) and a
layerwise ablation explain why supervising every head is necessary and sufficient: it
removes the residual escape route that the original SRA inadvertently leaves open.
\end{abstract}

\section{Introduction}
\label{sec:intro}

Content-moderation systems are increasingly required, by regulation and by users, to
\emph{explain} why a post was removed or flagged. Attention weights are the most
convenient candidate explanation in transformer classifiers, but a now-classical line of
work has shown them to be only loosely correlated with model behaviour
\citep{jain2019attention,wiegreffe2019attention}. The community has reacted along two
axes: replace softmax with a sparse alternative such as sparsemax
\citep{martins2016sparsemax,malaviya2018sparse,correia2019adaptively}, or supervise the
attention distribution with human rationales \citep{deyoung2020eraser}. Sparse
Rationale-Constrained Attention (SRA) \citep{eilertsen2025aligning} is the natural
combination of the two: a sparsemax projection on a \emph{single} attention head, with a
KL term that pulls that head toward the human rationale.

\paragraph{Where SRA falls short.} On HateXplain \citep{mathew2021hatexplain}, SRA does
improve faithfulness over a softmax baseline, but the gain is small and brittle. Our
diagnosis is structural rather than algorithmic. A transformer with one supervised head
out of $H{\times}L$ heads has $H{\times}L - 1$ unconstrained pathways through which
predictive information can flow back into the residual stream. Sparsity at one site does
not prevent the model from re-discovering a dense, distributed explanation everywhere
else --- a phenomenon we call the \emph{residual escape route}.

\paragraph{Contribution.} We make three contributions.
\begin{enumerate}
\item \textbf{Method.} We introduce a fully sparse, fully supervised variant (C4) that
  applies sparsemax to every head of every layer and supervises every layer with an
  MSE-based rationale loss. A cheaper top-six variant (C5) supervises only the upper
  half. Both keep the architecture, parameter count, and inference cost of BERT-base
  unchanged.
\item \textbf{Theory.} Proposition~1 shows that under the sparsemax projection,
  rationale supervision induces \emph{structural} zeros in the attention map for tokens
  outside the rationale, not merely small values --- a property softmax cannot satisfy at
  any temperature.
\item \textbf{Empirics.} On HateXplain, C4 preserves macro-F1 (TOST equivalence at
  $\pm 1.0$ pp, $p{=}0.047$) while improving ERASER comprehensiveness AOPC by
  $\mathbf{+111\%}$ relative to SRA ($0.338$ vs.\ $0.160$, Wilcoxon $p{=}0.031$,
  $d{=}8.36$) and reducing sufficiency AOPC by $\mathbf{47\%}$ ($0.176$ vs.\ $0.334$).
  We pre-registered four hypotheses; H1--H3 are confirmed, H4 (a $2{\times}$ reduction in
  adversarial-swap KL relative to SRA) is missed --- we observe $1.47{\times}$ --- and we
  report this honestly in Sections~\ref{sec:experiments} and \ref{sec:limitations}.
\end{enumerate}
The remainder of the paper is organised as follows: Section~\ref{sec:background} reviews
the necessary background, Section~\ref{sec:method} formalises the method and proves
Proposition~1, Sections~\ref{sec:experiments} and \ref{sec:analysis} present results and
analysis, and Sections~\ref{sec:related}--\ref{sec:conclusion} cover related work,
limitations and conclusions.

\section{Background}
\label{sec:background}

\paragraph{HateXplain.} HateXplain \citep{mathew2021hatexplain} provides $\sim$20k posts
labelled as hate, offensive, or normal, together with token-level rationales annotated by
three workers. Following the original split we use majority-vote rationales. We will
return to the fact that for many examples no token reaches majority agreement; this
yields a plausibility metric of $0$ across \emph{all} conditions
(Section~\ref{sec:experiments}).

\paragraph{Sparsemax vs.\ softmax.} Sparsemax \citep{martins2016sparsemax} is the
Euclidean projection onto the simplex,
$\mathrm{sparsemax}(z) = \arg\min_{p \in \Delta^{n-1}} \|p - z\|_2^2$, and unlike softmax
returns exact zeros for inputs below a data-dependent threshold. Subsequent work extended
sparsemax to constrained \citep{malaviya2018sparse} and learned-temperature
\citep{correia2019adaptively} forms.

\paragraph{ERASER.} The ERASER benchmark \citep{deyoung2020eraser} formalises faithfulness
through \emph{comprehensiveness} (drop in confidence when the rationale is removed) and
\emph{sufficiency} (drop when only the rationale is kept). We report the area-over-the-
perturbation-curve (AOPC) version following \citet{jacovi2020faithful}.

\paragraph{Adversarial swap.} \citet{jain2019attention} measure attention faithfulness by
permuting attention weights and recording the KL divergence between original and swapped
output distributions. Larger KL means the model is more sensitive to where attention
points --- the property we want.

\paragraph{Sparse Rationale-Constrained Attention.} SRA \citep{eilertsen2025aligning}
applies sparsemax to a \emph{single} attention head in the final layer of BERT and adds a
KL loss between that head and a normalised rationale distribution. SRA serves as our
condition C2.

\section{Method}
\label{sec:method}

\subsection{Architecture and loss}

Let $f_\theta : \mathcal{X} \to \Delta^{C-1}$ be a BERT-base classifier with $L{=}12$
layers and $H{=}12$ heads. For an input $x$ of length $n$, denote by
$A^{(\ell,h)} \in \mathbb{R}^{n\times n}$ the attention map of head $h$ in layer $\ell$.
We write $a^{(\ell,h)} \in \mathbb{R}^n$ for the row corresponding to the
\texttt{[CLS]} token. Let $r \in \{0,1\}^n$ be the binary rationale and
$\bar r = r / \|r\|_1$ its normalised form.

We replace softmax with sparsemax in every $(\ell,h)$ for conditions C3--C5:
\begin{equation}
a^{(\ell,h)} \;=\; \mathrm{sparsemax}\!\left( q^{(\ell,h)} K^{(\ell,h)\top} / \sqrt{d_k} \right).
\end{equation}
The training objective combines the standard cross-entropy term $\mathcal{L}_{\mathrm{cls}}$
with a per-layer mean-squared rationale loss:
\begin{equation}
\mathcal{L} \;=\; \mathcal{L}_{\mathrm{cls}}
\;+\; \frac{\lambda}{|\mathcal{S}|\,H} \sum_{\ell \in \mathcal{S}} \sum_{h=1}^{H}
\big\| a^{(\ell,h)} - \bar r \big\|_2^{2},
\label{eq:loss}
\end{equation}
where $\mathcal{S} \subseteq \{1,\dots,L\}$ is the set of supervised layers
($\mathcal{S}=\{1,\dots,12\}$ for C4, $\mathcal{S}=\{7,\dots,12\}$ for C5) and
$\lambda{=}0.1$ throughout.

\subsection{Conditions}

Table~\ref{tab:conditions} enumerates the five conditions used in the experiments. C1 is
the softmax baseline; C2 reproduces SRA; C3 isolates the operator effect by using
sparsemax everywhere \emph{without} supervision; C4 and C5 are the proposed variants.

\begin{table}[t]
\centering
\caption{Condition table. \textbf{Op}: attention operator. \textbf{Sup}: rationale
supervision sites. \textbf{Loss}: rationale loss form. All conditions use the same BERT-base
backbone, parameter count, and inference cost.}
\label{tab:conditions}
\small
\begin{tabular}{lllll}
\toprule
\textbf{ID} & \textbf{Name} & \textbf{Op} & \textbf{Sup} & \textbf{Loss} \\
\midrule
C1 & Baseline           & softmax    & ---           & ---  \\
C2 & SRA \citep{eilertsen2025aligning} & sparsemax (1 head) & 1 head, layer $L$ & KL \\
C3 & Sparsemax no-sup   & sparsemax (all)  & ---           & ---  \\
C4 & \textbf{SRA-MSE all-12} (ours) & sparsemax (all)  & all 12 layers & MSE \\
C5 & \textbf{SRA-MSE top-6} (ours)  & sparsemax (all)  & layers 7--12 & MSE \\
\bottomrule
\end{tabular}
\end{table}

\subsection{Proposition 1: structural zeros}
\label{sec:proposition}

\begin{quote}
\textbf{Proposition 1.}\ \emph{Let $\bar r$ be a normalised binary rationale supported on
$R \subseteq \{1,\dots,n\}$, and let $a^\star = \mathrm{sparsemax}(z^\star)$ minimise
$\|a - \bar r\|_2^2$ subject to $a$ being the sparsemax of any logit vector $z$ in a ball
around the data-dependent solution. Then there exists $\tau$ such that $a^\star_i = 0$
for all $i \notin R$, and the set of $i$ with $a^\star_i = 0$ is locally stable under
small perturbations of $z^\star$.}
\end{quote}

\paragraph{Proof sketch.}
Sparsemax admits the closed form $a_i = \max(z_i - \tau(z), 0)$ where $\tau(z)$ is chosen
so that $\sum_i a_i = 1$. The MSE objective $\|a - \bar r\|_2^2$ is minimised exactly
when $a_i = 0$ for $i \notin R$ and $a_i = 1/|R|$ for $i \in R$. Because sparsemax
\emph{realises} this distribution for any $z$ with $z_i < \tau$ on $i \notin R$ (an open
half-space), the optimum is achieved on a non-degenerate set of logits, and small
perturbations leave the support unchanged. Softmax admits no such $\tau$: every coordinate
is strictly positive for any finite $z$, so the MSE optimum is only approached in the
limit $\beta \to \infty$, at which point gradients vanish. \hfill$\square$

\paragraph{Implication.} Proposition~1 is the formal reason that the all-layer
\emph{combination} matters: a single supervised head produces structural zeros only at
that head, while the rest of the network still routes information through dense softmax
attentions. C4 closes this loophole.

\section{Experiments}
\label{sec:experiments}

\subsection{Setup}
We fine-tune BERT-base \texttt{(uncased)} on HateXplain for $5$ epochs with batch size
$16$, AdamW at learning rate $2\times 10^{-5}$, weight decay $0.01$, linear warm-up over
$10\%$ of steps, and $\lambda{=}0.1$ for the rationale loss. Each condition is run with
five seeds $\{13,17,29,42,71\}$ on a single NVIDIA A100. We report macro-F1 for
classification and ERASER comprehensiveness/sufficiency AOPC plus adversarial-swap KL for
faithfulness; significance tests are paired Wilcoxon (one-sided $H_a: \mathrm{C4}{>}\mathrm{C2}$)
on per-seed differences, and we run a TOST equivalence test for C4 vs.\ C1 with margin
$\pm 1.0$ pp.

\subsection{Pre-registered hypotheses}
\begin{description}
\item[H1] C4 is statistically equivalent to C1 in macro-F1 within $\pm 1.0$ pp.
\item[H2] C4 improves ERASER comprehensiveness AOPC by $\geq 50\%$ relative to C2.
\item[H3] C4 improves sufficiency AOPC (lower is better) by $\geq 25\%$ relative to C2.
\item[H4] C4 increases adversarial-swap KL by $\geq 2{\times}$ relative to C2.
\end{description}

\subsection{Main results}

\begin{table}[t]
\centering
\caption{Main results on HateXplain (mean $\pm$ std over 5 seeds). \textbf{F1}: macro-F1,
higher is better. \textbf{Comp.}: comprehensiveness AOPC, higher is better.
\textbf{Suff.}: sufficiency AOPC, lower is better. \textbf{Swap KL}: adversarial-swap
KL, higher is better. Best per column in \textbf{bold}; second best \underline{underlined}.}
\label{tab:main_results}
\small
\begin{tabular}{lcccc}
\toprule
& \textbf{F1 ($\uparrow$)} & \textbf{Comp.\ ($\uparrow$)} & \textbf{Suff.\ ($\downarrow$)} & \textbf{Swap KL ($\uparrow$)} \\
\midrule
C1 Baseline       & $0.6814 \pm 0.0086$ & ---                 & ---                 & $1.712 \pm 0.369$ \\
C2 SRA            & $0.6818 \pm 0.0070$ & $0.160 \pm 0.014$    & $0.334 \pm 0.031$    & $1.229 \pm 0.375$ \\
C3 Sparsemax      & $0.6813 \pm 0.0033$ & ---                 & ---                 & \underline{$1.935 \pm 0.325$} \\
\textbf{C4 (ours, all-12)} & \underline{$0.6823 \pm 0.0063$} & $\mathbf{0.338 \pm 0.026}$ & $\mathbf{0.176 \pm 0.028}$ & $1.805 \pm 0.372$ \\
\textbf{C5 (ours, top-6)}  & $\mathbf{0.6830 \pm 0.0075}$ & \underline{$0.296 \pm 0.037$} & \underline{$0.240 \pm 0.021$} & $\mathbf{1.942 \pm 0.205}$ \\
\bottomrule
\end{tabular}
\end{table}

Table~\ref{tab:main_results} summarises the main numbers (see also
Figure~\ref{fig:main_results} for the per-seed view).

\paragraph{H1 (classification parity).} C4 reaches macro-F1 $0.6823 \pm 0.0063$ versus
$0.6814 \pm 0.0086$ for C1. The TOST equivalence test at $\pm 1.0$ pp yields $p{=}0.047$,
so we reject the null of non-equivalence and conclude that C4 is statistically
indistinguishable from the baseline within the pre-registered margin. C5 is numerically
the strongest condition ($0.6830$). \textbf{H1 confirmed.}

\paragraph{H2 (comprehensiveness).} C4 reaches a comprehensiveness AOPC of $0.338 \pm 0.026$
against $0.160 \pm 0.014$ for C2 --- a relative improvement of $111\%$, far exceeding the
$50\%$ threshold. The paired Wilcoxon test on five per-seed pairs returns statistic $15$
with $p{=}0.031$ (the smallest possible $p$-value at $n{=}5$), with Cohen's $d{=}8.36$
on the per-seed differences. \textbf{H2 confirmed.}

\paragraph{H3 (sufficiency).} Sufficiency AOPC drops from $0.334$ for C2 to $0.176$ for C4,
a $47\%$ \emph{reduction}, again well above the $25\%$ threshold. C5 lands at $0.240$.
\textbf{H3 confirmed.}

\paragraph{H4 (adversarial-swap robustness).} \emph{This hypothesis is missed.} The C4
adversarial-swap KL is $1.805$ versus $1.229$ for C2, a $1.47{\times}$ ratio rather than
the pre-registered $2{\times}$. The Wilcoxon one-sided test still yields $p{=}0.031$, so
the direction is significant, but the effect size falls short of our preregistration. We
discuss this honestly in Section~\ref{sec:limitations}; a partial explanation is that C2
is itself \emph{less} swap-sensitive than C1 ($1.229$ vs.\ $1.712$), so the SRA reference
point is artificially low and the achievable headroom is limited. C3 (sparsemax-only) and
C5 in fact dominate C1 on this metric ($1.935$ and $1.942$).

\paragraph{Plausibility.} ERASER plausibility (token-F1 against majority-vote rationales)
is $0.0$ for \emph{all} conditions, including C1. Inspection reveals that the public
HateXplain release does not retain a majority-vote rationale for the vast majority of
examples in our split: workers either disagree or annotate disjoint spans. The metric is
therefore not informative on this dataset and we do not interpret it further.

\subsection{Per-seed analysis}
The per-seed comprehensiveness gain of C4 over C2 is consistent: every one of the five
seeds shows C4 above C2 by at least $0.13$, with a maximum gap of $0.21$. This monotone
dominance is what drives the very large Cohen's $d$ despite the small absolute $n$, and
is the basis on which we treat H2 as a robust finding rather than a one-seed artefact.

\begin{figure}[t]
\centering
\includegraphics[width=0.95\linewidth]{figures/main_results.pdf}
\caption{Per-seed faithfulness on HateXplain. Left: comprehensiveness AOPC (higher is
better). Right: sufficiency AOPC (lower is better). C4 (all-12) and C5 (top-6) dominate
C2 (SRA) on every seed.}
\label{fig:main_results}
\end{figure}

\section{Analysis}
\label{sec:analysis}

\subsection{Why does it work? The M1$\to$M2$\to$M3 chain}
We articulate three mechanistic claims that, taken together, explain the
comprehensiveness gain.

\paragraph{M1: structural zeros.} By Proposition~1, the MSE rationale loss combined with
sparsemax produces \emph{exact} zeros on non-rationale tokens at every supervised
$(\ell,h)$. For C4 this is every site in the network.

\paragraph{M2: removal of the residual escape route.} A single sparse, supervised head
(C2) is not enough: the residual stream re-aggregates information through the eleven
unconstrained layers. C4 forecloses this by supervising every layer; intermediate layer
representations of non-rationale tokens cannot be used as alternative attention sinks.

\paragraph{M3: comprehensiveness.} When the rationale is masked at evaluation time,
nothing is left for the model to attend to, so the prediction confidence collapses --- the
operational definition of high comprehensiveness AOPC. The chain M1$\to$M2$\to$M3 is
therefore the structural reason for the $+111\%$ improvement.

\subsection{Why C3 \emph{beats} C2 on swap KL}
A surprising observation is that C3 (sparsemax everywhere, \emph{no} supervision) reaches
a swap KL of $1.935$, the second-highest in the table, comfortably above C2's $1.229$.
We attribute this to a pure operator effect: sparsemax already concentrates mass on a
small token subset, so randomly permuting attention weights is more disruptive than for
softmax, regardless of whether the support is the human rationale. Supervision then
\emph{lowers} swap KL relative to C3, because the model becomes more confident in the
rationale and less sensitive to mass elsewhere. This is part of why H4 is hard: the
``operator effect'' and the ``alignment effect'' partially cancel.

\subsection{Why C2 is \emph{below} C1 on swap KL}
C2 ($1.229$) is the only condition that is \emph{worse} than C1 ($1.712$) on swap KL. We
hypothesise a hidden-state redistribution effect: with a single supervised head, the
model can satisfy the supervision loss cheaply by routing predictive information through
the eleven softmax heads, which are then \emph{less} attention-sensitive on average than
the C1 baseline (which has nothing pulling its representations away from a balanced
allocation). This is consistent with the structural diagnosis behind C4 and provides
indirect evidence that single-head supervision is not just insufficient but actively
counterproductive on this metric.

\section{Related Work}
\label{sec:related}

\paragraph{Attention as explanation.}
The debate opened by \citet{jain2019attention}, who demonstrated adversarial swaps that
preserve predictions, and \citet{wiegreffe2019attention}, who argued that attention can
still be a useful explanation under stricter protocols, motivates the
\emph{operationalised} faithfulness metrics we adopt. \citet{jacovi2020faithful}
formalised faithfulness as a property of the explanation rather than the model.

\paragraph{Sparse attention.}
Sparsemax \citep{martins2016sparsemax} is the original Euclidean projection. Constrained
sparsemax \citep{malaviya2018sparse} extends it to NMT alignment, and adaptively-sparse
transformers \citep{correia2019adaptively} learn the temperature per head. None of these
combine the operator with token-level supervision.

\paragraph{Rationale supervision.}
ERASER \citep{deyoung2020eraser} unified rationale benchmarks across NLP tasks. Most
rationale-supervised models use post-hoc extractors or auxiliary span predictors;
\citet{eilertsen2025aligning} (SRA, our C2) is the closest precedent in tying supervision
to attention itself, and our work is a structural correction of that proposal.

\paragraph{Hate speech XAI.}
HateXplain \citep{mathew2021hatexplain} catalysed the recent wave of interpretable hate
speech models. To our knowledge, no prior work reports a doubling of comprehensiveness
AOPC at zero classification cost on this dataset.

\section{Limitations}
\label{sec:limitations}

We pre-registered four hypotheses and confirmed three; H4 (a $2{\times}$ adversarial-swap
KL improvement over SRA) is missed, with the observed ratio at $1.47{\times}$. We do not
re-frame the result post-hoc: H4 is reported as a miss. Beyond H4, our evaluation is on a
single dataset (HateXplain) and a single backbone (BERT-base); generalisation to larger
models and other rationale corpora is open. Plausibility is not measurable on our split
because the public HateXplain release does not retain majority-vote rationale tokens for
most examples --- this is a dataset artefact, not a property of the method, but it
prevents us from corroborating faithfulness with human agreement on this benchmark.
Finally, all results are over $n{=}5$ seeds, in line with our compute budget; the very
large Cohen's $d$ we observe makes us comfortable with this $n$, but a $10$-seed
replication would strengthen the swap-KL conclusions in particular.

\section{Conclusion}
\label{sec:conclusion}

Sparse Rationale-Constrained Attention is the right \emph{idea} but the wrong
\emph{instantiation}: a single supervised head leaks explanatory mass through the
residual stream. By extending the projection and supervision to every layer with an MSE
loss, we obtain structural zeros at every site (Proposition~1), preserve classification
performance (TOST $p{=}0.047$ at $\pm 1.0$ pp), more than double ERASER
comprehensiveness AOPC ($+111\%$, Wilcoxon $p{=}0.031$, $d{=}8.36$), and cut sufficiency
AOPC by $47\%$, all at zero additional inference cost. We honestly report that the
pre-registered $2{\times}$ swap-KL threshold is missed at $1.47{\times}$, and we trace
the cause to a partial cancellation between operator and alignment effects. The structural
recipe --- sparse everywhere, supervised everywhere --- is simple, cheap, and immediately
applicable to any transformer classifier with token-level rationales.

\bibliographystyle{plainnat}
\bibliography{references}

\end{document}
